\documentclass[11pt]{article}
\usepackage[margin=1in]{geometry}
\usepackage[T1]{fontenc}
\usepackage{lmodern}
\usepackage{microtype}
\usepackage{setspace}
\usepackage{amsmath}
\usepackage{amssymb}
\usepackage{booktabs}
\usepackage{longtable}
\usepackage{array}
\usepackage{hyperref}
\usepackage{xcolor}

\hypersetup{colorlinks=true,linkcolor=blue!60!black,urlcolor=blue!60!black,citecolor=blue!60!black}
\setstretch{1.15}

\title{From Eight-Cell Design to Winner-Path Causal Chain:\\A Narrative Chapter on the Sparse Signal-Field SAE Project}
\author{Implementation Narrative Chapter}
\date{June 1, 2026}

\begin{document}
\maketitle

\begin{abstract}
This chapter reconstructs, in narrative form, what the Sparse Signal-Field SAE project set out to do, what it implemented, what it actually ran, where it diverged from its own original experimental matrix, and why those divergences happened. The project began with an explicit 2$\times$2$\times$2 design (E1--E8) that crossed training protocol, SAE family, and mask regime. Execution followed a different trajectory: first proving sparse reconstruction under conservative assumptions, then branching architecture and masking where uncertainty was highest, then focusing heavily on causal transition and chain behavior under manifold re-grounding. The final picture is not a failed plan, nor a completed factorial. It is a selective but scientifically coherent progression: strong evidence for sparse transform-domain sufficiency, strong evidence that raw multi-step code prediction is unstable, and strong evidence that re-grounding is the key mechanism that recovers compositional utility. At the same time, full planned Family II breadth was not completed, and that boundary is quantitatively stated here rather than softened. The objective of this chapter is to deliver a technically precise account that reads as a continuous argument, not a status memo.
\end{abstract}

\section{Introduction and Research Question}
The motivating question of this project is easiest to state by analogy. In natural image processing, raw pixel space is usually dense: most coefficients are nonzero, and directly imposing sparsity in that basis is poor compression. Yet transformed spaces such as wavelets and learned dictionaries reveal compressible structure. A small, strategically chosen subset of transformed coefficients can preserve most perceptually and statistically meaningful information. The project asks whether intermediate activations of a frozen vision model behave the same way.

If the answer is yes, then an activation field can be recoded into an overcomplete sparse basis, aggressively masked in that basis, decoded back to activation space, and passed through the frozen remainder of the model with minimal behavior change. If the answer is no, then sparse transform-domain language in this setting is either weak, layer-specific, or purely reconstruction-level without behavioral faithfulness.

The project therefore framed success in three nested levels. The first level was reconstruction sufficiency within each section: can sparse transformed coefficients rebuild hidden maps with small downstream damage. The second level was cross-section predictability: can sparse code at section $Q_t$ predict sparse code at section $Q_{t+1}$. The third level was chained depth-wise viability: can those relationships be composed from Q1 to Q5 without collapse.

This hierarchy of claims is important because each layer introduces a new failure mode. Reconstruction can succeed while transitions fail; transitions can succeed one step at a time while chains fail under repeated distribution shift. The project eventually encountered exactly this structure. In hindsight, that made the final results more informative than a single pass/fail criterion would have been.

\section{Original Design Intent: What E1--E8 Meant Conceptually}
The high-level specification in \texttt{high-level/sparse\_signal\_sae\_experiment\_plan\_high\_accuracy\_conv.tex} and the corresponding PDF proposes an explicit factorial grid with three binary dimensions. The first dimension is training protocol: Independent layerwise learning (Family I) versus Incremental dependent sparse-code learning (Family II). The second dimension is representation family: expressive convolutional Field SAE versus shared per-location Vector SAE. The third dimension is mask semantics: global TopK versus receptive-field-local TopK.

The eight planned experiments are therefore not arbitrary labels. They are a carefully chosen decomposition of causal uncertainty. E1 through E4 isolate architecture and mask behavior under independent layerwise code learning and then transition fitting across frozen code spaces. E5 through E8 ask whether the same architecture/mask choices still hold when code construction itself is trained with explicit dependence structure, rather than only connected later by predictors.

This distinction between protocol families is easy to erase in casual summaries and must be preserved in rigorous reporting. Family I can produce excellent reconstruction and even strong one-step transitions while still hiding chain instability under compounding off-manifold predictions. Family II was intended to attack that instability structurally by how dependent codes are learned, not merely by post hoc predictor regularization. Any claim that ``Family II was run'' must therefore specify whether it means full planned factorial dependent-code training or a narrower stabilization variant.

Another original design element that shaped execution is the expectation that both global and RF-local masking should be scientifically meaningful. Global masking answers whether sparse support can be concentrated anywhere in transformed space. RF-local masking answers a stricter question: whether each downstream receptive region can be supported by sparse local transformed coefficients. The second regime is theoretically closer to locality in hierarchical construction, but also algorithmically harsher.

The plan's structure, then, was conceptually sound: broad factorial coverage before downstream emphasis. The eventual divergence from this ordering came from empirical signals and practical constraints, not from a lack of implementation clarity.

\section{Experimental Substrate: Backbone, Sections, and Metrics}
Before discussing chronology, it is helpful to state what remained stable across phases. The backbone anchor for most reported results is frozen ResNet-56 on CIFAR-100 with 71.83\% top-1 (from \texttt{logs/backbone\_r56\_c100\_\_done.json}). The project taps five sections (Q1--Q5) and caches activations through \texttt{src/cache\_activations.py}, enabling repeatable SAE and transition experiments decoupled from full backbone retraining loops.

Each SAE maps activation $H\in\mathbb{R}^{C\times H\times W}$ to code $Z\in\mathbb{R}^{K\times H\times W}$ with overcompleteness $K=8C$ in the core runs. Masks are applied in code space, not activation space. Reconstruction then decodes masked code $\tilde{Z}$ back to $\hat{H}$, and the frozen tail computes downstream logits. This exact splice-and-tail setup is central: it turns reconstruction quality into behavioral preservation metrics rather than only geometric distortion metrics.

Reported metrics therefore come in families. Reconstruction metrics include relative $\ell_2$ and cosine similarity in normalized section space. Downstream metrics include spliced top-1 and KL divergence versus original logits. Sparsity metrics include effective retained fraction, which is especially informative when comparing global fraction targeting to RF-local fixed-k schemes where implicit retained fractions differ strongly by resolution.

Controls were intentionally built into the substrate. Random-dictionary runs and PCA baselines test whether observed retention-performance behavior is an artifact of dimensionality or a learned transform-domain effect. Without those controls, strong SAE numbers would be less informative.

\section{What Was Actually Built: Architecture, Training, and Workflow Evolution}
The implemented codebase is not a thin prototype around one experiment; it is a modular system with explicit interfaces for reconstruction, transition, chain evaluation, intervention, and diagnostics. The implementation details matter because they determine what the eventual results can and cannot imply.

In \texttt{src/sae.py}, Field SAE and Vector SAE are materially distinct. Field SAE introduces local spatial processing through residual blocks with depthwise kernels and pointwise expansions. Vector SAE uses shared per-location linear mappings (1x1 convolutions) in a more standard sparse coding style with decoder normalization. Both maintain the design constraint that sparse code is not bypassed by direct skip from input activation to output activation.

In \texttt{src/masks.py}, global TopK and RF-local TopK are implemented with deterministic hard selection. Global masking retains a target fraction over all $(k,h,w)$ coordinates. RF-local masking partitions or approximates receptive windows and retains local top coefficients within each region. This implementation allows direct experimental contrast between globally concentrated sparse support and locality-constrained sparse support.

In \texttt{src/train\_sae.py}, layerwise training uses staged sparsification: warmup without mask pressure, annealing from generous retention to target budget, then fixed target retention. This is a practical stability choice that emerged as necessary for robust convergence under hard TopK constraints.

In \texttt{src/transitions.py} and \texttt{src/train\_transition.py}, the project formalized Family I transition learning as predictor fitting between frozen source and destination SAE spaces. But the same module also defines \texttt{frozen\_block\_hybrid}, which turned out to be one of the most consequential engineering additions. By decoding sparse source code and propagating through the real frozen backbone block before re-encoding at destination, hybrid mode cleanly separates representational sufficiency from predictor error.

In \texttt{src/eval\_chain.py}, chain evaluation became explicitly multimodal: raw prediction chaining, remasking chained predictions, full decode-reencode remanifold projection (re-grounding), and hybrid chain. This was not present in the original matrix abstraction but became necessary once chain collapse appeared.

In \texttt{src/train\_chain.py}, a further shift occurred: scheduled sampling and optional BPTT through re-grounding. Predictors are no longer trained only in teacher-forced local conditions; they are exposed to self-generated carriers, with or without gradient through the manifold projection path.

At the workflow layer, \texttt{scripts/run\_grid.py} reveals the strategic decision structure. Phase 1 is E3 anchor only. Phase 2 is E1/E2/E4 branch. There is no built-in E5--E8 full matrix scheduler. Later scripts target transition and tier campaigns. This is where planned breadth gave way to winner-path depth.

\section{What Was Actually Run: The Chronological Story}
\subsection{Phase 1 as Feasibility Anchor Rather Than Competitive Sweep}
The first major campaign, represented in \texttt{runs/phase1}, was intentionally narrow. It did not attempt to win the architecture contest. It attempted to settle the base claim: does sparse transform-domain reconstruction preserve model behavior at all. Running E3 first (Vector+Global under Independent protocol) made this gate conservative. If even this baseline failed, broader branching would be premature.

It did not fail. Across sections and budgets, the model retained substantial downstream fidelity at modest retention fractions, with stronger retention in deeper sections under aggressive sparsity than naive intuition might predict. Equally important, random-dictionary controls in \texttt{runs/phase1\_rand} were near chance on deeper sections, and PCA baselines in \texttt{logs/baselines\_pca.json} trailed learned SAE behavior under comparable or larger effective budgets. The implication was straightforward: success was not due to trivial projection or latent high-dimensional redundancy alone; learned sparse basis structure mattered.

From a planning perspective, this phase changed the epistemic status of the core hypothesis from speculative to operational. It also justified spending compute on architecture and mask branching rather than debating feasibility in the abstract.

\subsection{Phase 2 as the Architectural and Masking Pivot}
Phase 2 executed E1, E2, and E4 across all five sections (30 cells total). This branch had two goals: identify whether expressive spatial mixing materially improves low-budget performance, and test whether RF-local masking is usable or too brittle.

The first goal resolved clearly in early sections. At 2\% retained under global masking, Field+Global outperformed Vector+Global by large margins at Q1 and Q2 and a meaningful margin at Q3. The second goal produced an even sharper contrast. Under RF-local masking at \texttt{k\_rf=4}, Field remained viable while Vector degraded substantially, especially at intermediate resolutions where local support constraints were strict and the per-location encoder had insufficient local context mixing.

This phase did not merely produce better numbers; it produced a direction. It identified a robust winner path for downstream hierarchy experiments: Field architecture with global 5\% as stable primary operating point, with RF-local retained as a sparse but less uniform option.

\subsection{Phase 3: Strong Single-Step Edges, Failed Raw Chain}
With winner SAEs selected, Phase 3 trained adjacent predictors (Q1$\rightarrow$Q2, Q2$\rightarrow$Q3, Q3$\rightarrow$Q4, Q4$\rightarrow$Q5) in \texttt{runs/phase3/T\_*}. Single-step outcomes were partly excellent. The first edge nearly matched hybrid upper bound. Later edges remained useful, with Q3$\rightarrow$Q4 clearly the hardest segment.

The conceptual break came at composition. \texttt{runs/phase3/chain\_result.json} shows raw chained predictions collapsing to 0.0116, despite a hybrid chain at 0.7122 and original at 0.7183. This pattern rules out one simplistic explanation and confirms another. It rules out ``sparse support has insufficient information'' because hybrid succeeds. It confirms ``predictor composition drifts off manifold under recursive exposure'' because raw chain fails while single-step edges are nontrivial.

This is the point where the project's center of gravity moved from ``which architecture wins'' to ``what mechanism keeps code-space dynamics composable.''

\subsection{Phase 3b: Re-grounding as Mechanism, Not Patch}
Phase 3b in \texttt{runs/phase3b} and \texttt{runs/phase3b\_bptt} trained transitions with scheduled sampling and re-grounded carriers, still anchored to frozen winner SAEs. Re-grounded chain performance improved from 0.5930 (phase3 eval mode) to 0.6624 with scheduled sampling, then to 0.7010 when training through re-grounding with BPTT. Hybrid remained 0.7122.

This trajectory matters because it demonstrates a graded mechanism. If re-grounding were a brittle post-processing trick, one would expect little trainable gain around it. Instead, gains were substantial and monotonic under stronger training alignment with chain-time conditions. Conversely, raw chain remained near chance even after these training upgrades, indicating that unprojected dense predictor outputs remain poor carriers for repeated composition.

At this stage, the project had an internally consistent narrative: sparse support is sufficient, local transitions can be learned, but chain stability depends on manifold projection at each step. This was a stronger explanatory framework than simply reporting a failed chain and moving on.

\subsection{Tier Diagnostics: Why the Mechanism Interpretation Held}
Tier diagnostics tested whether the chain interpretation aligned with causal and representation-quality evidence.

Intervention results in \texttt{runs/intervene\_result.json} showed top-ablation damage ratios above random baselines by factors such as 3.29x and 7.5x, with additional locality concentration beyond area-null expectations. This supports that learned sparse coordinates are not arbitrary compression artifacts; they encode concentrated causal influence.

Feature audits in \texttt{runs/audit/audit\_Q3.json} and \texttt{runs/audit/audit\_Q5.json} showed low dead-feature rates and limited duplication, reducing concern that gains were driven by degenerate dictionaries. The stronger Q3$\rightarrow$Q4 predictor stress test improved over chance but did not close the main gap, reinforcing that the hard edge was not solved by simply deepening transition network capacity.

Together, these diagnostics supported the mechanistic pivot rather than contradicting it.

\subsection{Phase 4 Cluster Reruns: Scale, Robustness, and Transfer After the Blocked Node}
The later cluster reruns matter because they answer a different class of question than Phases 1--3b. The earlier phases established the core mechanism on the main backbone. The reruns ask whether that mechanism generalizes sensibly when one changes dataset/backbone scale, repeats the winner under multiple seeds, or moves to adjacent transfer settings.

The R20/C10 grid is best read as a scale-down and dataset-shift stress test of the winner recipe. It reruns Field+Global-style sparse reconstruction on ResNet-20/CIFAR-10 across Q1--Q5 and fractions 2\%, 5\%, and 10\%. All 15 cells completed. The important pattern is not merely that the jobs finished, but that the deep sections were extremely strong and largely saturated by 5\% retained. Q3, Q4, and Q5 all sat around 0.92 spliced top-1 across the tested fractions. In other words, the sparse-sufficiency story does not look delicately tuned to the ResNet-56/CIFAR-100 anchor.

The seed grid answers a narrower but important criticism: perhaps the winner-path claims depend on one lucky training seed. Here the project reran the 5\%-retained winner configuration across seeds 0, 1, and 2 for every section. All 15 cells completed. The quantitative spread is small, especially in the most relevant later sections. Q3 ranged from 0.7108 to 0.7142, Q4 from 0.7151 to 0.7184, and Q5 from 0.7165 to 0.7174. This does not prove full statistical robustness in a formal sense, but it does materially strengthen confidence that the headline results are not seed accidents.

The taxonomy experiment asks a more conceptual question about what sort of representation best carries the Q4$\rightarrow$Q5 relationship. Sparse mode outperformed dense and cross controls, with 0.7063 spliced top-1 versus 0.6902 and 0.6099. That result is important because it supports a stronger claim than ``sparse codes can reconstruct.'' It suggests that the sparse carrier is itself an especially useful transition representation in this same-shape section pair.

The ResNet-110 reconstruction sweep extends the winner recipe to a deeper backbone. All five sections completed successfully, with spliced top-1 values of 0.7316, 0.7311, 0.7226, 0.7290, and 0.7318 from Q1 to Q5. The familiar Q3 weakness persists, which is scientifically useful. It means the project's hardest section is not an isolated fluke of one backbone, but a recurrent structural difficulty near the same representational regime change.

Finally, the optional ViT transfer run moved from ``implemented but blocked'' to ``executed with a real artifact.'' Its metrics are not yet the centerpiece of the project, but they are nontrivial enough to matter: prediction agreement 0.9440 and KL 0.0427 at 5\% retained. For the narrative, the main significance is that the project's transfer story now includes a concrete output rather than only an aspiration.

\section{Results in Integrated Interpretation}
A memo-style report can list many metrics. A chapter should articulate what those metrics jointly imply.

At the reconstruction level, evidence is robust. Field SAE under global masking preserved near-baseline behavior with aggressive sparsity in representative layers, and controls established that this was learned-structure-dependent, not random or linear projection trivia.

At the architecture/mask interaction level, results indicate that local expressive mixing before sparse bottlenecking is not optional when moving from global to RF-local constraints. The Vector regime can remain competitive under global conditions at easier depths but degrades sharply when local sparsity budgets become strict.

At the transition level, one-step learned predictors are often strong, yet this strength does not extrapolate under naive composition. This disconnect is the central lesson of Phase 3. Multi-step failure was not due to absent information, as hybrid chain remained strong, but due to representational carrier mismatch across recursive predictor use.

At the chain-stability level, re-grounding transformed the narrative from binary failure to conditional success. Chain-aware training and BPTT-through-projection recovered most of the hybrid gap. The residual gap likely reflects both persistent hard-edge mismatch (especially around Q3$\rightarrow$Q4 geometry changes) and remaining mismatch between predictor families and true inter-section dynamics.

In short, the project confirmed sparse sufficiency and revealed a precise compositional constraint: sparse-code chains are usable when repeatedly projected back onto destination SAE manifolds, and unstable otherwise.

\section{Clarification of Family I and Family II Execution}
This section resolves the central ambiguity that generated disagreement in prior summaries.

Family I, in both plan and execution, means independently trained SAEs with learned transitions between frozen code spaces. E1--E4 were executed as explicit matrix cells and have direct artifact coverage in \texttt{runs/phase1} and \texttt{runs/phase2}.

Family II in the design plan means incremental dependent sparse-code learning across the same SAE/mask matrix breadth. In that strict sense, full E5--E8 were not executed as matrix-complete cells.

However, Family II-style execution did occur in phase3b form. The system trained transitions under scheduled sampling and re-grounded carriers, optionally with BPTT through projection, on the winner configuration. This is not ``none'' and not ``full.'' It is partial and specific.

Quantitatively, configuration breadth is one realized Family II-style configuration out of four planned Family II configuration cells, approximately 25\% coverage of planned Family II matrix breadth. Within that one configuration, chain depth coverage is complete across all four adjacent transitions from Q1 to Q5.

This framing is the only one that is both technically accurate and rhetorically fair. Saying Family II ``was not run'' discards real phase3b artifacts. Saying Family II ``was run'' without qualifiers overstates matrix completion. The correct statement is: Family II-style mechanism work was run deeply on one winner path; full Family II matrix breadth remains incomplete.

\section{Limitations, Blocked Runs, and Pending Work}
Not all incompleteness in this project reflects scientific indecision. Some of it reflects systems instability documented in run logs.

Tier logs in \texttt{runs/tier\_rest} do show repeated CUDA initialization and dependency failures that initially prevented completion of taxonomy, ViT transfer, and several downstream extension stages. But those failure logs are only the first half of the story. A later cluster rerun wave completed the same extension set successfully: \texttt{runs/r20c10/grid\_summary.json} now shows 15 successful R20/C10 cells, \texttt{runs/seeds/grid\_summary.json} shows 15 successful seed-grid cells, \texttt{runs/taxonomy/taxonomy\_q4q5.json} is present, \texttt{runs/r110/F\_Q1} through \texttt{runs/r110/F\_Q5} are populated, and \texttt{runs/vit\_sae/vit\_result.json} exists after later dependency resolution. The correct historical statement is therefore not that these experiments never ran. It is that they were genuinely blocked on the first node and later recovered on cluster reruns.

Scientifically, the main limitation is the mismatch between planned and executed Family II breadth. The project has strong winner-path chain stabilization evidence but not the full cross-architecture, cross-mask dependent-protocol comparison originally envisioned in E5--E8.

A second limitation is that robustness and generality are now quantified unevenly rather than absent outright. The successful seed grid improved robustness accounting for the winner setup, and the successful taxonomy, ResNet-110, and ViT reruns improved generality accounting beyond the core backbone. Even so, these extensions remain thinner and less replicated than the main ResNet-56 winner-path evidence.

These limitations do not erase the positive findings, but they constrain the scope of defensible claims. The strongest claims are about the verified winner path and verified mechanisms, not about all potential configurations in the original matrix.

\section{Concluding Synthesis}
The most accurate way to summarize this project is not ``plan completed'' and not ``plan abandoned.'' The project executed a scientifically coherent narrowing strategy under empirical and systems constraints.

It began with a conservative anchor and demonstrated sparse reconstruction sufficiency against meaningful controls. It then identified decisive architecture-mask interactions that selected a winner path. It discovered a critical mismatch between single-step transition quality and naive chain composition. It introduced and validated a mechanism, re-grounding, that explains and repairs most of that mismatch when integrated into training. It supported this mechanism with intervention and feature-health diagnostics rather than relying on one aggregate metric.

What remains is equally clear. Full planned Family II matrix breadth was not executed; only one of four planned Family II configuration cells received deep chain-aware treatment. The chapter therefore closes with a dual conclusion: the sparse transform-domain hypothesis and re-grounded chain mechanism are strongly evidenced on the executed path, and the full dependent-protocol factorial remains pending work rather than completed work.

That conclusion is not a compromise between optimism and caution. It is the direct consequence of the artifact record.

\section{Mechanistic Deep Dive: Why Re-grounding Changed the Outcome}
The phrase ``re-grounding'' can sound procedural until one asks what exact mismatch it repairs. In the project, each transition predictor is trained to map one sparse code distribution to the next under teacher-like inputs that are close to section-specific SAE manifolds. During free composition, however, each predictor receives the previous predictor's output, and this output is generally dense, slightly mis-scaled, and not guaranteed to obey the destination SAE geometry that training implicitly assumed.

From the perspective of dynamical systems, raw chaining compounds model mismatch. Let $F_t$ denote predictor $t\\rightarrow t+1$, and let $\\mathcal{M}_{t+1}$ denote the destination SAE manifold after encode-plus-mask. Teacher-forced training reduces local error around $\\mathcal{M}_t$, but raw rollout repeatedly applies $F_t$ to points that drift away from $\\mathcal{M}_t$. Even if each one-step predictor has acceptable local approximation error, the composition can diverge rapidly when each next input comes from an out-of-distribution region generated by the previous approximation.

The hybrid chain result is decisive in this context because it proves the model is not fundamentally information-starved at the sparsity budget used. In hybrid mode, sparse support at one section is decoded, pushed through the true frozen block, then re-encoded for the next section. This procedure keeps each step on a valid representational path. The observed 0.7122 chain top-1 near a 0.7183 original baseline therefore shows that sparse support can, in principle, sustain end-to-end behavior through depth.

Raw chain collapse to 0.0116 does not contradict hybrid success; it isolates predictor-carrier incompatibility. The project then tested two increasingly strong interventions that can be interpreted as projections toward $\\mathcal{M}_{t+1}$. Re-masking alone imposes sparsity shape but does not enforce realistic decoded structure, and performance rose only to 0.2126. Full re-grounding, decode then re-encode then mask, enforces both representational shape and destination encoder compatibility, and performance rose to 0.5930.

Phase 3b's scheduled-sampling regime and BPTT-through-re-grounding then addressed a second-order issue: even if evaluation projects onto manifold each step, predictors trained mostly on teacher carriers still underperform when exposed to self-generated carriers. Scheduled sampling increases exposure during training, improving re-grounded chain to 0.6624. BPTT through the projection path aligns gradients with chain-time behavior, pushing re-grounded chain to 0.7010.

This progression provides a coherent causal interpretation: the dominant bottleneck in multi-step sparse-code chaining is not lack of sparse information and not merely insufficient predictor capacity; it is recursive off-manifold drift. Re-grounding is effective because it directly targets that drift at each composition boundary.

The Q3$\\rightarrow$Q4 edge still remains harder than others. Here the model undergoes a resolution transition and representation regime change, and small errors in how code maps across this boundary can produce larger downstream distortion than equal-sized errors on resolution-preserving edges. The stronger Q3$\\rightarrow$Q4 predictor experiment improved over chance but remained below the best baseline edge behavior, reinforcing that this bottleneck reflects structural mapping difficulty rather than a simple underparameterized block.

If one abstracts this outcome beyond this project, the lesson is portable: sparse-code transitions are not ordinary dense latent dynamics. They are constrained by section-specific geometry induced by encoder-decoder pairs and hard masking operators. Any chain architecture that ignores this geometry can look acceptable under one-step diagnostics yet fail catastrophically in rollout. Any architecture that acknowledges this geometry can recover substantial compositional performance even without changing the underlying frozen backbone.

\section{Protocol Counterfactual: What Full E5--E8 Would Have Required}
Because the project explicitly planned E5--E8, it is useful to state what ``full completion'' would have looked like, not as blame assignment but as technical contrast. In strict planned Family II terms, each of the four architecture-mask combinations would require dependent sparse-code training procedures that make later section codes genuinely conditioned by earlier sparse representations during code construction, not only by downstream transition fitting after independent SAE spaces are fixed.

Concretely, this would imply protocol-specific runs for Field+Global, Field+RF-local, Vector+Global, and Vector+RF-local under dependent-code learning assumptions, all with corresponding reconstruction and downstream preservation checks, and with chain or adjacent-transition evaluations comparable enough to support cross-cell interpretation. It would also imply storing result artifacts in clearly separated E5--E8 directories that mirror E1--E4 reporting style.

What happened instead was a targeted compression of this space. After Phase 2 identified a clear winner path for practical hierarchy work, the project invested in mechanism development on that path: frozen winner SAEs, scheduled sampling, re-grounded carriers, and optional BPTT through projection. This produced high-value insight but only on one configuration breadth-wise.

The quantitative coverage statement is therefore straightforward. Planned Family II breadth had four configuration cells. Executed Family II-style breadth had one winner-path configuration. That is approximately 25\\% breadth coverage. Within that one breadth cell, depth coverage across four adjacent transitions was complete in the sense of full Q1$\\rightarrow$Q5 chain handling.

This asymmetry has interpretation consequences. It permits strong claims about mechanism viability on the selected configuration, moderate claims about generality to nearby configurations that already showed strong Family I behavior, and weak claims about combinations not executed under Family II-style training. In particular, it leaves open whether RF-local Family II for Field would preserve the same stabilization gains and whether Vector-family dependent-protocol variants would exhibit similar or worse chain sensitivity.

The counterfactual also clarifies why disagreements emerged in plain-language summaries. A reader focused on matrix breadth can reasonably say ``E5--E8 were not run.'' A reader focused on protocol mechanics can reasonably say ``Family II-style training was run.'' Both statements are true under different axes. The correct synthesis requires naming the axis explicitly: breadth incomplete, mechanism substantial on one path.

From a planning standpoint, this chapter's correction is not merely semantic hygiene. It prevents a common failure in technical communication: flattening protocol-level nuance into binary completion language. Future readers deciding what to run next need the exact map of what is known, where it is known, and why that evidence may or may not transfer.

\section{Extended Chronology of Decisions, Evidence, and Tradeoffs}
A continuous narrative benefits from reconstructing not just results but decision points. The project had several such points where available evidence materially changed what was rational to run next.

The first decision point followed Phase 1 anchor completion. At that point, the project had enough evidence to reject the null that sparse coding was merely cosmetic in this setting. Trained SAE behavior differed sharply from random and linear controls, and downstream preservation at practical budgets was credible. The rational next step was not to deepen one baseline forever, but to interrogate architecture and mask interactions where uncertainty remained high.

The second decision point followed Phase 2 branching. The Field family was not marginally better; it was materially more robust at aggressive budgets in early sections and notably more resilient under RF-local constraints. Once this pattern emerged across sections, spending scarce experimental budget equally across all architecture variants became less defensible for the hierarchy question. Winner-path focus became a deliberate opportunity-cost decision.

The third decision point followed Phase 3 chain collapse under raw composition. Without hybrid evidence, this could have been interpreted as ``sparse hierarchy hypothesis failed.'' With hybrid evidence, that interpretation became untenable. Hybrid near-baseline chain showed that sparse support carried sufficient information. Thus the immediate question changed from ``is hierarchy false'' to ``what transformation restores compositional stability of learned predictors.''

The fourth decision point followed early re-grounding improvements. A modest gain from remasking suggested that sparsity-shape alignment helps but is insufficient. A larger gain from full decode-reencode projection suggested manifold alignment is the major lever. This justified moving beyond evaluation-time patching to training-time alignment via scheduled sampling and BPTT-through-projection.

The fifth decision point concerned whether capacity scaling alone could solve hard edges, especially Q3$\\rightarrow$Q4. The stronger predictor run improved but did not eliminate the gap relative to easier edges and hybrid references. This reduced confidence in ``just scale predictor'' as the primary strategy and strengthened confidence in representation-geometry explanations.

The sixth decision point was operational rather than conceptual: tier campaign interruption under node/GPU instability. At first, some pending runs were blocked for infrastructure reasons despite implemented scripts and codepaths. The project then responded by packaging those runs for cluster relaunch, which changed the evidence base. The current tree therefore no longer ends with pure blockage. It ends with uneven completeness of a different kind: deep evidence on selected paths, completed but thinner extension evidence on scale, robustness, and transfer runs, and explicit log history showing both the failures and the later successful reruns.

Viewed end to end, this chronology is not random drift from plan. It is a sequence of evidence-conditioned reallocations: from feasibility to architecture discrimination, from architecture discrimination to causal-transition diagnostics, from diagnostics to mechanistic stabilization, and from mechanism to later generalization attempts that were first blocked and then partially recovered on cluster reruns.

This matters for interpreting ``why not full E5--E8 first.'' In a perfect environment with unlimited compute and no runtime interruptions, broad matrix completion before mechanism deepening might have been preferable for symmetry. In this environment, once strong winner and failure signatures appeared, mechanism-deepening offered higher marginal information per run than symmetric completion of weak candidates. That choice generated a stronger causal explanation but a weaker breadth-completeness profile.

One might worry that this strategy introduces confirmation bias by over-investing in a winner. That is a legitimate concern in general, but here two factors mitigate it. First, the winner was selected from direct cross-cell evidence in Phase 2 rather than anecdotal preference. Second, mechanism tests included failure-preserving controls (raw chain remains near chance) and hard-edge stress tests (Q3$\\rightarrow$Q4 remains difficult), which reduced the risk of narrating only positive outcomes.

Another concern is whether re-grounding ``cheats'' by injecting information unavailable to raw prediction. In this project's formulation, re-grounding decodes and re-encodes predicted representations through fixed SAEs and fixed masks. It does not access future labels or hidden states from outside the chain. It enforces representational consistency with the learned sparse manifold. Whether one labels this as architectural prior or projection operator, it is a legitimate mechanism within the model class being tested, and its causal role is empirically measurable through side-by-side chain variants.

The final tradeoff is reporting style itself. A matrix-centric report emphasizes completion percentages. A mechanism-centric report emphasizes explanatory depth on executed paths. This chapter intentionally combines both: it preserves exact completion boundaries while explaining why the executed subset still yields strong scientific signal.

\section{Interpretive Boundaries and What Can Be Claimed Responsibly}
Given the evidence, the strongest responsible claim is conditional and layered. The project demonstrates that frozen CNN activations are highly compressible in learned sparse transform-domain coordinates while preserving downstream behavior at low retained fractions. It further demonstrates that sparse code at one section can support useful prediction of sparse code at the next, and that end-to-end chain viability is recoverable when each step is re-aligned to destination SAE manifold structure.

A stronger claim---that all planned dependent-protocol variants would show similar behavior---is not yet warranted. The artifact record does not contain that breadth. A weaker claim---that chain behavior is hopeless under sparse coding---is also contradicted by re-grounded and hybrid evidence. The truth sits between these extremes: sparse-code hierarchy is viable under geometric discipline and incomplete in breadth coverage.

This boundary-focused interpretation is not rhetorical caution for its own sake. It is what allows future work to be targeted rather than repetitive. If the dominant unknown is protocol breadth generalization, then the next high-value runs are not another winner-path retread but controlled Family II-style executions on remaining architecture-mask cells, ideally with shared reporting templates that permit direct comparability to existing phase3b artifacts.

Finally, communicating this boundary clearly protects downstream users of these results. Engineers reading this chapter should not conclude that E5--E8 are done. Reviewers should not conclude Family II is absent. Both would be factual errors. The chapter's job is to prevent exactly those errors by keeping mechanism depth and matrix breadth as separate accounting dimensions.

\appendix
\section{Appendix A: Compact Status Matrix}
\begin{table}[h!]
\centering
\caption{Planned E1--E8 versus run coverage in current artifacts.}
\begin{tabular}{llll}
\toprule
ID & Protocol & SAE/Mask & Status \\
\midrule
E1 & Family I (Independent) & Field + Global & Executed (10 cells) \\
E2 & Family I (Independent) & Field + RF-local & Executed (10 cells) \\
E3 & Family I (Independent) & Vector + Global & Executed (15 cells) \\
E4 & Family I (Independent) & Vector + RF-local & Executed (10 cells) \\
E5 & Family II (Planned Incremental) & Field + Global & Not executed as full grid cell \\
E6 & Family II (Planned Incremental) & Field + RF-local & Not executed as full grid cell \\
E7 & Family II (Planned Incremental) & Vector + Global & Not executed as full grid cell \\
E8 & Family II (Planned Incremental) & Vector + RF-local & Not executed as full grid cell \\
\bottomrule
\end{tabular}
\end{table}

\section{Appendix B: Key Numeric Anchors}
\begin{longtable}{p{0.33\linewidth}p{0.61\linewidth}}
\toprule
Metric family & Representative values from run artifacts \\
\midrule
Backbone reference & ResNet-56/CIFAR-100: 71.83\% (\texttt{logs/backbone\_r56\_c100\_\_done.json}) \\
E1 vs E3 at 2\% (Q1-Q5) & Q1 0.7139 vs 0.6376; Q2 0.7164 vs 0.6327; Q3 0.7078 vs 0.6646; Q4 0.7106 vs 0.7098; Q5 0.7147 vs 0.7123 \\
E2 vs E4 (k\_rf=4) & Q1 0.7158 vs 0.6858; Q2 0.6992 vs 0.4645; Q3 0.6698 vs 0.2430; Q4 0.6949 vs 0.5942; Q5 0.6443 vs 0.2504 \\
Phase 3 transitions & Q1$\rightarrow$Q2 0.7146 (hybrid 0.7151); Q2$\rightarrow$Q3 0.6907 (hybrid 0.7105); Q3$\rightarrow$Q4 0.6504 (hybrid 0.7113); Q4$\rightarrow$Q5 0.7089 (hybrid 0.7149) \\
Chain metrics (Phase 3 eval) & raw 0.0116; re-mask 0.2126; re-ground 0.5930; hybrid 0.7122; original 0.7183 \\
Phase 3b & scheduled-sampling re-grounded chain 0.6624; BPTT re-grounded chain 0.7010; hybrid 0.7122 \\
N8 intervention & top/random importance ratios: 3.29x, very large at Q2$\rightarrow$Q3 due near-zero random effect, 4.00x, 7.50x \\
Feature audit & Q3 dead 2.73\%, dup 5.86\%; Q5 dead 1.37\%, dup 0.00\% \\
\bottomrule
\end{longtable}

\section{Appendix C: Reproducibility Artifact Map}
Primary execution and evidence paths referenced in this chapter:
\begin{itemize}
    \item Planning source: \texttt{high-level/sparse\_signal\_sae\_experiment\_plan\_high\_accuracy\_conv.tex} and PDF
    \item Grid orchestrator: \texttt{scripts/run\_grid.py}
    \item Transition training: \texttt{src/train\_transition.py}
    \item Chain training/eval: \texttt{src/train\_chain.py}, \texttt{src/eval\_chain.py}
    \item Core runs: \texttt{runs/phase1}, \texttt{runs/phase2}, \texttt{runs/phase3}, \texttt{runs/phase3b}, \texttt{runs/phase3b\_bptt}
    \item Causal and quality diagnostics: \texttt{runs/intervene\_result.json}, \texttt{runs/audit/audit\_Q3.json}, \texttt{runs/audit/audit\_Q5.json}
    \item Baseline controls: \texttt{logs/baselines\_pca.json}, \texttt{runs/phase1\_rand}
    \item Blocked-tier logs: \texttt{runs/tier\_rest/*}, \texttt{runs/r20c10/grid\_summary.json}, \texttt{runs/seeds/grid\_summary.json}
\end{itemize}

\end{document}
